\documentclass[11pt,a4paper]{article}

% ============================================================
%  Guía de implementación HU9 / HU11 / HU12 / HU14 — Nanys Care
%  Tema adecuado a la paleta de la app (rosa / lila / naranja)
%  Compilar con: pdflatex (2 pasadas)
% ============================================================

\usepackage[T1]{fontenc}
\usepackage[utf8]{inputenc}
\usepackage[spanish,es-noindentfirst]{babel}
\usepackage{lmodern}
\usepackage{amssymb}
\usepackage[scaled=0.95]{helvet}
\renewcommand{\familydefault}{\sfdefault}

\usepackage[a4paper,margin=2.4cm,headheight=15pt]{geometry}
\usepackage[table]{xcolor}
\usepackage{graphicx}
\usepackage{enumitem}
\usepackage{booktabs}
\usepackage{tabularx}
\usepackage{array}
\usepackage{titlesec}
\usepackage{fancyhdr}
\usepackage[most]{tcolorbox}
\usepackage{listings}
\usepackage{microtype}
\usepackage[hidelinks]{hyperref}

% ---------- Paleta del tema (AppColors) ----------
\definecolor{appPrimary}{HTML}{E91E63}   % Rosa cálido
\definecolor{appPrimaryDark}{HTML}{C2185B}
\definecolor{appPrimaryLight}{HTML}{F8BBD0}
\definecolor{appSecondary}{HTML}{7E57C2} % Lila
\definecolor{appAccent}{HTML}{FFB74D}    % Naranja suave
\definecolor{appSuccess}{HTML}{4CAF50}
\definecolor{appWarning}{HTML}{FFA726}
\definecolor{appError}{HTML}{E53935}
\definecolor{appInfo}{HTML}{29B6F6}
\definecolor{appText}{HTML}{2D2D2D}
\definecolor{appTextSec}{HTML}{6B6B6B}
\definecolor{appBg}{HTML}{FFF8F9}
\definecolor{codeBg}{HTML}{2D2233}
\definecolor{codeKw}{HTML}{F48FB1}
\definecolor{codeStr}{HTML}{B39DDB}
\definecolor{codeCom}{HTML}{9E9E9E}

\hypersetup{
  colorlinks=true,
  linkcolor=appPrimaryDark,
  urlcolor=appSecondary,
  pdftitle={Guia de implementacion HU9 HU11 HU12 HU14 - Nanys Care}
}

% ---------- Encabezado / pie ----------
\pagestyle{fancy}
\fancyhf{}
\renewcommand{\headrulewidth}{0pt}
\fancyhead[L]{\footnotesize\color{appTextSec}\textbf{Nanys\,Care}}
\fancyhead[R]{\footnotesize\color{appTextSec}Guía de implementación}
\fancyfoot[C]{\footnotesize\color{appTextSec}\thepage}

% ---------- Títulos de sección ----------
\titleformat{\section}
  {\Large\bfseries\color{appPrimaryDark}}
  {}{0pt}
  {\colorbox{appPrimaryLight}{\makebox[0.55em]{}}\hspace{0.6em}}
\titlespacing*{\section}{0pt}{1.6em}{0.8em}

\titleformat{\subsection}
  {\large\bfseries\color{appSecondary}}{}{0pt}{}
\titlespacing*{\subsection}{0pt}{1.1em}{0.4em}

% ---------- Estilo de código (Dart) ----------
\lstdefinelanguage{Dart}{
  keywords={abstract,async,await,class,const,extends,final,for,if,return,
    void,Future,List,bool,double,int,String,var,new,this,required,super,
    continue,true,false,null,import,not},
  morekeywords=[2]{Widget,BuildContext,StatelessWidget,FutureBuilder,Column,
    Text,Card,ListTile,Icon,Scaffold,AppBar},
  sensitive=true,
  morecomment=[l]{//},
  morecomment=[s]{/*}{*/},
  morestring=[b]',
  morestring=[b]"
}
\lstset{
  language=Dart,
  backgroundcolor=\color{codeBg},
  basicstyle=\ttfamily\scriptsize\color{white},
  keywordstyle=\color{codeKw}\bfseries,
  keywordstyle=[2]\color{appInfo},
  stringstyle=\color{codeStr},
  commentstyle=\color{codeCom}\itshape,
  numbers=left,
  numberstyle=\tiny\color{codeCom},
  numbersep=8pt,
  showstringspaces=false,
  breaklines=true,
  frame=single,
  framesep=8pt,
  rulecolor=\color{codeBg},
  xleftmargin=6pt,
  aboveskip=10pt,
  belowskip=10pt,
  columns=fullflexible,
  keepspaces=true,
}

% ---------- Cajas (callouts) ----------
\newtcolorbox{nota}[1][Nota]{
  colback=appInfo!8, colframe=appInfo, coltitle=white,
  fonttitle=\bfseries\small, title=#1, boxrule=0.8pt, arc=4pt,
  left=8pt,right=8pt,top=6pt,bottom=6pt}
\newtcolorbox{tip}[1][Recomendación]{
  colback=appSecondary!8, colframe=appSecondary, coltitle=white,
  fonttitle=\bfseries\small, title=#1, boxrule=0.8pt, arc=4pt,
  left=8pt,right=8pt,top=6pt,bottom=6pt}
\newtcolorbox{listo}[1][Hecho cuando...]{
  colback=appSuccess!8, colframe=appSuccess, coltitle=white,
  fonttitle=\bfseries\small, title=#1, boxrule=0.8pt, arc=4pt,
  left=8pt,right=8pt,top=6pt,bottom=6pt}

\newcommand{\hu}[1]{\textcolor{appPrimary}{\textbf{#1}}}
\setlist[itemize]{leftmargin=1.4em,itemsep=2pt,topsep=3pt}
\setlist[enumerate]{leftmargin=1.6em,itemsep=3pt,topsep=3pt}

% ============================================================
\begin{document}

% ---------- Portada ----------
\begin{titlepage}
\thispagestyle{empty}
\begin{tcolorbox}[colback=appPrimary,colframe=appPrimary,arc=10pt,
  left=24pt,right=24pt,top=28pt,bottom=28pt,halign=left]
  {\color{white}\Huge\bfseries Guía de implementación}\\[6pt]
  {\color{appPrimaryLight}\LARGE HU9 · HU11 · HU12 · HU14}
\end{tcolorbox}

\vspace{1.4cm}
{\color{appText}\Large\bfseries Nanys Care}\\[2pt]
{\color{appTextSec}\large Aplicación de cuidado infantil — Flutter}

\vspace{0.8cm}
{\color{appTextSec}
Plan técnico para implementar las historias de usuario de\\
\textbf{recordatorios automáticos}, \textbf{calificación de cuidadores},
\textbf{evaluación de tutores} y \textbf{recordatorios de citas},
respetando la arquitectura actual del proyecto
(modelos + servicios \texttt{ChangeNotifier} + \texttt{StorageService}).}

\vfill
\begin{tcolorbox}[colback=appBg,colframe=appPrimaryLight,arc=8pt,
  left=16pt,right=16pt,top=12pt,bottom=12pt]
\renewcommand{\arraystretch}{1.4}
\begin{tabularx}{\linewidth}{>{\bfseries\color{appPrimaryDark}}l X}
HU9  & Recordatorios automáticos \\
HU11 & Calificar cuidadores \\
HU12 & Evaluar tutores \\
HU14 & Recordatorios de citas \\
\end{tabularx}
\end{tcolorbox}

\vspace{0.6cm}
{\color{appTextSec}\small Documento técnico generado para el proyecto
\texttt{blasphatrocraft\_sa\_nannyscare} (rama \texttt{version1}).}
\end{titlepage}

% ---------- Índice ----------
\renewcommand{\contentsname}{\textcolor{appPrimaryDark}{Contenido}}
\tableofcontents
\vspace{1em}

% ============================================================
\section{Estado actual y alcance}

La guía cubre cuatro historias de usuario. El proyecto ya tiene la base
(modelos \texttt{Resena} y \texttt{Notificacion}, almacenamiento local y
pantallas por rol), por lo que en varios casos solo hay que completar el flujo.

\begin{center}
\renewcommand{\arraystretch}{1.5}
\begin{tabularx}{\linewidth}{c >{\raggedright\arraybackslash}p{4.5cm} >{\raggedright\arraybackslash}X}
\toprule
\rowcolor{appPrimaryLight}
\textbf{HU} & \textbf{Título} & \textbf{Estado actual} \\
\midrule
\hu{HU9}  & Recordatorios automáticos & No existe el motor de recordatorios. \\
\hu{HU11} & Calificar cuidadores      & Parcial: ya se crea la \texttt{Resena}; falta promedio, evitar duplicados y mostrarlas. \\
\hu{HU12} & Evaluar tutores           & Parcial: ya existe la ``nota privada''; falta pantalla para consultarlas. \\
\hu{HU14} & Recordatorios de citas    & Depende de HU9. \\
\bottomrule
\end{tabularx}
\end{center}

\begin{nota}[HU9 + HU14 van juntas]
HU9 es el \textbf{motor} (servicio que genera recordatorios) y HU14 es su
\textbf{primer uso} (recordatorios de citas próximas).
\end{nota}

\subsection*{Orden recomendado}
\begin{enumerate}
  \item \hu{HU11} — completa lo ya empezado; es la más rápida.
  \item \hu{HU12} — reutiliza casi todo de HU11.
  \item \hu{HU9 + HU14} — el motor nuevo de recordatorios.
\end{enumerate}

% ============================================================
\section{HU11 — Calificar cuidadores}

Hoy \texttt{tutor\_citas\_screen.dart} (\texttt{\_calificar}) ya guarda una
\texttt{Resena}, pero \emph{no} actualiza el promedio del cuidador, permite
calificar la misma cita varias veces y las reseñas no se muestran.

\subsection{Paso 1 — Crear un \texttt{ResenaService}}
Centraliza la lógica para no repetir código entre tutor y cuidador.
Crea \texttt{lib/services/resena\_service.dart}:

\begin{lstlisting}
import 'package:flutter/foundation.dart';
import '../models/resena.dart';
import 'storage_service.dart';

class ResenaService extends ChangeNotifier {
  final StorageService _storage;
  ResenaService(this._storage);

  // Resenas PUBLICAS recibidas por un usuario (para mostrar en su perfil).
  Future<List<Resena>> getResenasPublicas(String destinatarioId) async {
    final all = await _storage.getResenas();
    final list = all
        .where((r) => r.destinatarioId == destinatarioId && !r.esPrivada)
        .toList();
    list.sort((a, b) => b.createdAt.compareTo(a.createdAt));
    return list;
  }

  // Notas PRIVADAS escritas por un autor (HU12 / RF15).
  Future<List<Resena>> getNotasPrivadas(String autorId) async {
    final all = await _storage.getResenas();
    final list = all.where((r) => r.autorId == autorId && r.esPrivada).toList();
    list.sort((a, b) => b.createdAt.compareTo(a.createdAt));
    return list;
  }

  // Evita duplicados: el autor ya califico esta cita?
  Future<bool> yaCalifico(String citaId, String autorId) async {
    final all = await _storage.getResenas();
    return all.any((r) => r.citaId == citaId && r.autorId == autorId);
  }

  Future<void> crearResena(Resena resena) async {
    await _storage.addResena(resena);
    if (!resena.esPrivada) {
      await _storage.recalcularCalificacionCuidador(resena.destinatarioId);
    }
    notifyListeners();
  }
}
\end{lstlisting}

\subsection{Paso 2 — Recalcular el promedio en \texttt{StorageService}}
Añade este método en \texttt{lib/services/storage\_service.dart}
(sección RESEÑAS):

\begin{lstlisting}
// Recalcula y persiste calificacionPromedio segun resenas publicas.
Future<void> recalcularCalificacionCuidador(String cuidadorId) async {
  final resenas = await getResenas();
  final publicas = resenas
      .where((r) => r.destinatarioId == cuidadorId && !r.esPrivada)
      .toList();
  if (publicas.isEmpty) return;

  final promedio =
      publicas.fold<double>(0, (s, r) => s + r.calificacion) / publicas.length;

  final perfiles = await getPerfilesCuidador();
  final idx = perfiles.indexWhere((p) => p.userId == cuidadorId);
  if (idx < 0) return;

  final p = perfiles[idx];
  perfiles[idx] = PerfilCuidador(
    userId: p.userId, descripcion: p.descripcion,
    nivelExperiencia: p.nivelExperiencia, certificaciones: p.certificaciones,
    capacidades: p.capacidades, tarifaPorHora: p.tarifaPorHora,
    ubicacion: p.ubicacion, diasDisponibles: p.diasDisponibles,
    horarioDisponible: p.horarioDisponible,
    calificacionPromedio: double.parse(promedio.toStringAsFixed(1)),
    totalServicios: p.totalServicios, verificado: p.verificado,
  );
  await savePerfilesCuidador(perfiles);
}
\end{lstlisting}

\begin{tip}
Más limpio: añade un \texttt{copyWith} a \texttt{PerfilCuidador} (como ya lo
tiene \texttt{Cita}) y úsalo aquí. Recomendado, opcional.
\end{tip}

\subsection{Paso 3 — Registrar el servicio en \texttt{main.dart}}
\begin{lstlisting}
final resenas = ResenaService(storage);
// ...y dentro del MultiProvider:
ChangeNotifierProvider<ResenaService>.value(value: resenas),
\end{lstlisting}
Recuerda agregar el campo y el \texttt{required this.resenas} en el constructor
de \texttt{NanysCareApp}, igual que los demás servicios.

\subsection{Paso 4 — Usar el servicio en \texttt{tutor\_citas\_screen.dart}}
En \texttt{\_calificar}, valida duplicados antes del diálogo y al enviar usa
\texttt{crearResena} en vez de \texttt{storage.addResena}:

\begin{lstlisting}
final resenaSvc = context.read<ResenaService>();
if (await resenaSvc.yaCalifico(cita.id, auth.currentUser!.id)) {
  // SnackBar "Ya calificaste esta cita" y return
}
// ...al confirmar:
await resenaSvc.crearResena(resena); // recalcula el promedio
\end{lstlisting}

\subsection{Paso 5 — Mostrar las reseñas en el perfil del cuidador}
En \texttt{cuidador\_detalle\_screen.dart}, agrega una sección ``Reseñas''
(es \texttt{StatelessWidget}, usa \texttt{FutureBuilder}):

\begin{lstlisting}
_SectionTitle('Resenas'),
FutureBuilder<List<Resena>>(
  future: context.read<ResenaService>().getResenasPublicas(user.id),
  builder: (context, snap) {
    if (!snap.hasData) return const SizedBox.shrink();
    final resenas = snap.data!;
    if (resenas.isEmpty) return const Text('Aun no tiene resenas');
    return Column(
      children: resenas.map((r) => ListTile(
        leading: const Icon(Icons.star, color: AppColors.warning),
        title: Text('${r.autorNombre}  ${r.calificacion.toStringAsFixed(1)}'),
        subtitle: Text(r.comentario),
      )).toList(),
    );
  },
),
\end{lstlisting}

\begin{listo}
El tutor califica una cita completada, \emph{no} puede repetir, el promedio del
cuidador sube y las reseñas aparecen en su perfil.
\end{listo}

% ============================================================
\section{HU12 — Evaluar tutores}

\texttt{cuidador\_solicitudes\_screen.dart} (\texttt{\_pedirNotaPrivada}) ya crea
una \texttt{Resena} con \texttt{esPrivada: true} al completar la cita. Falta una
pantalla donde el cuidador \textbf{consulte} esas evaluaciones.

\subsection{Decisión de diseño}
\begin{itemize}
  \item \textbf{Privada (RF15):} la evaluación del tutor solo la ve quien la
        escribió. Es lo que ya hace el código; mantenlo así.
  \item Si tu HU12 exige que \textbf{administrador/supervisor} vean las
        evaluaciones de tutores, entonces no la marques privada o añade un rol
        autorizado a leerlas.
\end{itemize}

\subsection{Paso 1 — Reusar \texttt{ResenaService}}
Cambia \texttt{storage.addResena(...)} por:
\begin{lstlisting}
await context.read<ResenaService>().crearResena(resena); // esPrivada: true
\end{lstlisting}
y valida duplicados con \texttt{yaCalifico(cita.id, autorId)}, igual que en HU11.

\subsection{Paso 2 — Pantalla ``Mis notas de tutores''}
Crea \texttt{lib/screens/cuidador/cuidador\_notas\_screen.dart}:

\begin{lstlisting}
class CuidadorNotasScreen extends StatelessWidget {
  const CuidadorNotasScreen({super.key});

  @override
  Widget build(BuildContext context) {
    final user = context.watch<AuthService>().currentUser!;
    return Scaffold(
      appBar: AppBar(title: const Text('Mis notas de tutores')),
      body: FutureBuilder<List<Resena>>(
        future: context.read<ResenaService>().getNotasPrivadas(user.id),
        builder: (context, snap) {
          if (!snap.hasData) {
            return const Center(child: CircularProgressIndicator());
          }
          final notas = snap.data!;
          if (notas.isEmpty) {
            return const Center(child: Text('Aun no tienes notas privadas'));
          }
          return ListView(
            padding: const EdgeInsets.all(16),
            children: notas.map((n) => Card(
              child: ListTile(
                leading: const Icon(Icons.lock_outline),
                title: Text('${n.calificacion.toStringAsFixed(1)} estrellas'),
                subtitle: Text(n.comentario.isEmpty ? '(sin comentario)' : n.comentario),
              ),
            )).toList(),
          );
        },
      ),
    );
  }
}
\end{lstlisting}

\subsection{Paso 3 — Enlazar la pantalla}
Agrega un acceso desde \texttt{cuidador\_home.dart} o
\texttt{cuidador\_perfil\_screen.dart} (un \texttt{ListTile}/botón con
\texttt{Navigator.push} a \texttt{CuidadorNotasScreen}).

\begin{listo}
Al completar una cita el cuidador evalúa al tutor y puede volver a ver esas
evaluaciones en ``Mis notas de tutores''.
\end{listo}

% ============================================================
\section{HU9 + HU14 — Recordatorios automáticos / de citas}

No hay backend ni \emph{cron}, así que el patrón realista es: \textbf{revisar las
citas al arrancar la app o al refrescar el home y generar notificaciones de tipo
\texttt{recordatorio} para las citas próximas que aún no tengan una.}

\subsection{Paso 1 — Crear \texttt{ReminderService}}
Crea \texttt{lib/services/reminder\_service.dart}:

\begin{lstlisting}
import '../models/cita.dart';
import '../models/notificacion.dart';
import 'storage_service.dart';
import 'notification_service.dart';

// HU9/HU14: genera recordatorios automaticos de citas proximas.
class ReminderService {
  final StorageService _storage;
  final NotificationService _notifications;
  ReminderService(this._storage, this._notifications);

  Future<void> generarRecordatoriosPendientes() async {
    final ahora = DateTime.now();
    final citas = await _storage.getCitas();
    final existentes = await _storage.getNotificaciones();

    for (final c in citas) {
      if (c.estado != EstadoCita.aceptada) continue;   // solo confirmadas
      final faltan = c.fechaInicio.difference(ahora);
      if (faltan.isNegative) continue;                  // ya paso
      if (faltan.inHours > 24) continue;                // dentro de 24h

      bool yaTiene(String userId) => existentes.any((n) =>
          n.tipo == TipoNotificacion.recordatorio &&
          n.usuarioId == userId && n.mensaje.contains(c.id));

      final cuando = _fmt(c.fechaInicio);
      if (!yaTiene(c.tutorId)) {
        await _notifications.enviarNotificacion(
          usuarioId: c.tutorId, titulo: 'Recordatorio de cita',
          mensaje: 'Tu cita con ${c.cuidadorNombre} es el $cuando. [${c.id}]',
          tipo: TipoNotificacion.recordatorio);
      }
      if (!yaTiene(c.cuidadorId)) {
        await _notifications.enviarNotificacion(
          usuarioId: c.cuidadorId, titulo: 'Recordatorio de cita',
          mensaje: 'Tienes una cita con ${c.tutorNombre} el $cuando. [${c.id}]',
          tipo: TipoNotificacion.recordatorio);
      }
    }
  }

  String _fmt(DateTime d) =>
      '${d.day}/${d.month}/${d.year} ${d.hour}:${d.minute}';
}
\end{lstlisting}

\begin{nota}[Detección de duplicados]
El \texttt{[\$\{c.id\}]} dentro del mensaje permite detectar duplicados sin tocar
el modelo. Si prefieres algo más limpio, agrega un campo opcional \texttt{citaId}
a \texttt{Notificacion} (actualizando \texttt{toJson}/\texttt{fromJson}).
\end{nota}

\subsection{Paso 2 — Registrar y disparar}
En \texttt{main.dart}:
\begin{lstlisting}
final reminders = ReminderService(storage, notifications);
await reminders.generarRecordatoriosPendientes(); // al arrancar
\end{lstlisting}
Vuelve a llamarlo al entrar al home o en el \emph{pull-to-refresh} de
\texttt{notificaciones\_screen.dart} (dentro de \texttt{\_load}):
\begin{lstlisting}
await context.read<ReminderService>().generarRecordatoriosPendientes();
\end{lstlisting}

\subsection{Paso 3 (opcional) — HU9 más amplio}
El mismo servicio puede generar otros recordatorios:
\begin{itemize}
  \item recordar al tutor \textbf{calificar} una cita completada sin reseña (enlaza con HU11);
  \item recordar al cuidador \textbf{completar} una cita aceptada cuya \texttt{fechaFin} ya pasó.
\end{itemize}
La UI ya está lista: \texttt{notificaciones\_screen.dart} muestra el ícono de
alarma y color naranja para \texttt{TipoNotificacion.recordatorio}.

\begin{listo}
Al abrir la app, las citas aceptadas dentro de las próximas 24\,h generan una
notificación de recordatorio (sin duplicarse) visible para tutor y cuidador.
\end{listo}

% ============================================================
\section{Checklist y pruebas}

\subsection*{Checklist}
\begin{itemize}[label=$\square$]
  \item \texttt{ResenaService} creado y registrado en \texttt{main.dart}.
  \item \texttt{StorageService.recalcularCalificacionCuidador} añadido.
  \item HU11: validación de duplicado + promedio actualizado + reseñas visibles.
  \item HU12: \texttt{crearResena} con \texttt{esPrivada:true} + pantalla ``Mis notas''.
  \item \texttt{ReminderService} creado y registrado.
  \item HU9/HU14: recordatorios al arrancar y al refrescar notificaciones.
  \item \texttt{flutter analyze} sin errores y \texttt{flutter run} probado.
\end{itemize}

\subsection*{Cómo probar rápido (datos demo)}
\begin{center}
\renewcommand{\arraystretch}{1.4}
\begin{tabularx}{\linewidth}{>{\bfseries\color{appPrimaryDark}}l X}
Tutor     & \texttt{carlos@nanyscare.com} / \texttt{123456} \\
Cuidadora & \texttt{maria@nanyscare.com} / \texttt{123456} \\
\end{tabularx}
\end{center}

\textbf{Flujo:} el tutor agenda cita $\rightarrow$ la cuidadora la acepta
$\rightarrow$ la cuidadora la ``marca completada'' (evalúa al tutor, HU12)
$\rightarrow$ el tutor califica (HU11) $\rightarrow$ reabrir la app para ver los
recordatorios (HU9/HU14). Para reiniciar datos: \texttt{StorageService.reset()}.

\vspace{1.2em}
\begin{center}
\color{appTextSec}\small\rule{0.4\linewidth}{0.4pt}\\[4pt]
Nanys Care · Documento técnico interno
\end{center}

\end{document}
